\setcounterpageref{ex}{0}
\TaoPhieuKiemTra{03}

\setcounter{ex}{0}
\setcounter{bt}{0}

\textbf{\large I. PHẦN TRẮC NGHIỆM (3,0 điểm)}

{\itshape Khoanh tròn vào chữ cái đặt trước câu trả lời đúng. Mỗi câu đúng được 0,5 điểm.}

\begin{ex} % Câu 1
Số lớn nhất có hai chữ số là
\choice
{\True $99$}
{$90$}
{$100$}
{$98$}
\loigiai{
Số lớn nhất có hai chữ số là $99$.\\
\textbf{Đáp án đúng: A}
}
\end{ex}

\begin{ex} % Câu 2
Trong phép tính $74 - 31 = 43$, số $31$ được gọi là
\choice
{Số bị trừ}
{\True Số trừ}
{Hiệu}
{Tổng}
\loigiai{
Số đứng sau dấu trừ gọi là Số trừ.\\
\textbf{Đáp án đúng: B}
}
\end{ex}

\begin{ex} % Câu 3
Độ dài $1\text{ dm}$ bằng bao nhiêu xăng-ti-mét?
\choice
{$1\text{ cm}$}
{$100\text{ cm}$}
{\True $10\text{ cm}$}
{$5\text{ cm}$}
\loigiai{
$1\text{ dm} = 10\text{ cm}$.\\
\textbf{Đáp án đúng: C}
}
\end{ex}

\begin{ex} % Câu 4
Ba điểm cùng nằm trên một đường thẳng được gọi là
\choice
{Ba điểm song song}
{Ba điểm cắt nhau}
{Ba đỉnh tam giác}
{\True Ba điểm thẳng hàng}
\loigiai{
Khi cùng nằm trên một đường thẳng, ba điểm đó thẳng hàng.\\
\textbf{Đáp án đúng: D}
}
\end{ex}

\begin{ex} % Câu 5
So sánh: $8 + 7 \dots 9 + 5$.
\choice
{$<$}
{\True $>$}
{$=$}
{Không so sánh được}
\loigiai{
$8 + 7 = 15$; $9 + 5 = 14$. Vì $15 > 14$ nên điền dấu $>$.\\
\textbf{Đáp án đúng: B}
}
\end{ex}

\begin{ex} % Câu 6
Tính: $8 + 6 - 5 = \dots$
\choice
{\True $9$}
{$10$}
{$8$}
{$7$}
\loigiai{
$8 + 6 = 14$; $14 - 5 = 9$.\\
\textbf{Đáp án đúng: A}
}
\end{ex}

\vspace{4mm}
\textbf{\large II. PHẦN TỰ LUẬN (7,0 điểm)}

{\itshape Thí sinh trình bày chi tiết bài làm vào giấy kiểm tra.}

\begin{bt}[1,0 điểm]
Tính nhẩm:
\begin{multicols}{2}
\begin{enumerate}[a)]
  \item $9 + 7$
  \item $7 + 8$
  \item $15 - 8$
  \item $17 - 9$
\end{enumerate}
\end{multicols}
\loigiai{
\begin{itemize}
  \item $9 + 7 = 16$
  \item $7 + 8 = 15$
  \item $15 - 8 = 7$
  \item $17 - 9 = 8$
\end{itemize}
}
\end{bt}

\begin{bt}[1,5 điểm]
Đặt tính rồi tính:
\begin{multicols}{3}
\begin{enumerate}[a)]
  \item $29 + 41$
  \item $81 - 30$
  \item $47 + 50$
\end{enumerate}
\end{multicols}
\loigiai{
\begin{itemize}
  \item $29 + 41 = 70$
  \item $81 - 30 = 51$
  \item $47 + 50 = 97$
\end{itemize}
}
\end{bt}

\begin{bt}[1,0 điểm]
Tính giá trị biểu thức kèm đơn vị đo:
\begin{enumerate}[a)]
  \item $19\text{ kg} + 10\text{ kg} - 9\text{ kg}$
  \item $28\text{ }l - 14\text{ }l + 9\text{ }l$
\end{enumerate}
\loigiai{
\begin{enumerate}[a)]
  \item $19\text{ kg} + 10\text{ kg} - 9\text{ kg} = 29\text{ kg} - 9\text{ kg} = 20\text{ kg}$.
  \item $28\text{ }l - 14\text{ }l + 9\text{ }l = 14\text{ }l + 9\text{ }l = 23\text{ }l$.
\end{enumerate}
}
\end{bt}

\begin{bt}[1,5 điểm]
Bạn An có $18$ bức tranh. Bạn Bình có nhiều hơn bạn An $7$ bức tranh. Hỏi bạn Bình có bao nhiêu bức tranh?
\loigiai{
Số bức tranh bạn Bình có là:
\[ 18 + 7 = 25 \text{ (bức tranh)} \]
\textbf{Đáp số:} $25$ bức tranh.
}
\end{bt}

\begin{bt}[1,0 điểm]
Một thùng chứa đầy $38\text{ }\text{kg}$ đường kính. Buổi sáng cửa hàng đã bán ra $15\text{ }\text{kg}$. Hỏi trong thùng còn lại bao nhiêu ki-lô-gam đường kính?
\loigiai{
Trong thùng còn lại số ki-lô-gam đường kính là:
\[ 38 - 15 = 23\text{ }\text{kg} \]
\textbf{Đáp số:} $23\text{ }\text{kg}$ đường kính.
}
\end{bt}

\begin{bt}[1,0 điểm]
Tìm số bé nhất có hai chữ số mà tổng hai chữ số của nó bằng $9$.
\loigiai{
Để số bé nhất thì chữ số hàng chục bé nhất, chọn chữ số hàng chục là $1$.\\
Chữ số hàng đơn vị: $9 - 1 = 8$.\\
Vậy số đó là $18$.\\
\textbf{Đáp số:} $18$.
}
\end{bt}
